% ============================================================
% 最小可编译 LaTeX 论文模板（贴合国赛真实优秀论文惯例，v2.0）
% 用途：避免编译失败重试浪费 token，从此模板开始撰写
% 编译：xelatex paper.tex （两次）
%
% 结构依据：组委会官方展示优秀论文的主流结构：
%   摘要页(标题+摘要+关键词，无目录、无队名日期)
%   → 问题重述 → 问题分析 → 模型假设 → 符号说明
%   → 模型的建立与求解（按问题组织，推导内联正文）
%   → 灵敏度与误差分析 → 模型评价 → 参考文献 → 附录
% ============================================================

\documentclass[12pt,a4paper]{ctexart}

% ---- 宏包 ----
\usepackage{amsmath,amssymb,amsthm}
\usepackage{graphicx}
\usepackage{booktabs}
\usepackage{geometry}
\usepackage{hyperref}
\usepackage{xcolor}
\usepackage{listings}
\usepackage{float}
\usepackage{caption}
\usepackage{subcaption}
\usepackage{enumitem}

% ---- 中文代码注释保留配置（历年校准反馈强化）----
% 1. 推荐使用 \lstinputlisting{problemX.py} 直接引用源文件，可避免中文注释丢失
% 2. 若需内联代码段，必须配置 extendedchars=true + escapeinside 以支持 CJK 字符
% 3. 等宽中文字体配置（可选）：\setCJKmonofont{SimSun} 或 Microsoft YaHei Mono
% 详见 .trae/papers/国赛公开惯例.md 二.1 "附录代码中文处理强化"

% ---- 页面设置（贴近真实国赛 Word 论文：宽边距 + 大页边空白）----
\geometry{left=3.0cm, right=3.0cm, top=2.54cm, bottom=2.54cm}

% ============================================================
% 颜值硬约束（与真实国赛论文逐页视觉比对后固化，禁止删改）
% 对标样本：真实国赛获奖论文
% ============================================================
% 1. 页码页脚居中、全文无页眉（ctexart 默认 headings 会产生页眉+右上角页码，
%    与真实论文"页脚居中页码、无页眉"不符，必须改 plain）
\pagestyle{plain}
% 2. 行距 1.4 倍（真实论文为 Word 1.5 倍/固定值 22 磅的宽松观感，
%    LaTeX 默认行距明显偏紧，是"一眼不像国赛论文"的主因之一）
\linespread{1.4}\selectfont
% 3. 标题体系（真实论文：一级标题"一、 xxx"黑体三号居中；
%    二级标题"1.1 xxx"黑体四号左对齐；三级黑体小四左对齐）
\ctexset{
    section = {
        format = \heiti\zihao{3}\centering,
        name = {,、},
        number = \chinese{section},
        aftername = \quad
    },
    subsection = {
        format = \heiti\zihao{4},
        number = \arabic{section}.\arabic{subsection},
        aftername = \quad
    },
    subsubsection = {
        format = \heiti\zihao{-4},
        number = \arabic{section}.\arabic{subsection}.\arabic{subsubsection},
        aftername = \quad
    }
}
% 4. 表格线粗（真实论文三线表顶/底线约 1.5 磅粗线，booktabs 默认偏细）
\setlength{\heavyrulewidth}{1.5pt}
\setlength{\lightrulewidth}{0.75pt}
% 5. 表格列对齐：符号说明等三线表建议列居中（真实论文为居中风格），
%    即 {ccc} 而非 {lll}；表格整体 \small 亦可接受
% ============================================================

% ---- 超链接设置 ----
\hypersetup{
    colorlinks=true,
    linkcolor=black,
    citecolor=black,
    urlcolor=blue
}

% ---- 代码 listings 设置 ----
% extendedchars=true + escapeinside 用于保留中文注释（详见上方 A-05/C-01 配置说明）
\lstset{
    basicstyle=\ttfamily\small,
    breaklines=true,
    showstringspaces=false,
    frame=single,
    numbers=left,
    numberstyle=\tiny\color{gray},
    keywordstyle=\color{blue},
    commentstyle=\color{green!50!black},
    stringstyle=\color{red},
    language=Python,
    tabsize=4,
    extendedchars=true,
    escapeinside={(*@}{@*)}
}

% ---- 图表路径 ----
\graphicspath{{figures/}}

\begin{document}

% ============================================================
% 摘要页（官方硬约束：≤1 页；只含标题+摘要+关键词）
% 禁止：\maketitle（会带入队名/日期）、\tableofcontents（官方明确不要目录）
% ============================================================
\begin{center}
    {\zihao{3}\heiti 论文标题（由 paper-writer 填写，准确简练醒目）}
\end{center}

\bigskip

\noindent\textbf{摘要：}

[第 1 段：用 1-3 句交代问题背景与总体思路。]

[以下按问题逐段：针对问题X，建立了什么模型（方法名），用什么工具/算法求解，
得到了什么结论——结论必须定量化（具体数值、百分比、最优方案参数）。
每问一段；若某问含两个子任务，可在段内分述。]

\bigskip

\noindent\textbf{关键词：}[关键词1；关键词2；关键词3；关键词4（3-5 个，含核心方法与核心指标）]

\newpage

% ============================================================
% 一、问题重述
% ============================================================
\section{问题重述}

\subsection{问题背景}

[问题背景描述。若题目含专有名词/附件说明，可在此处用三线表给出"名词解释与附件说明"。]

\subsection{具体问题重述}

\textbf{问题1}：[用自己的语言提炼为明确的数学任务，避免照搬原文]

\textbf{问题2}：[同上]

\textbf{问题3}：[同上]

\textbf{问题4}：[同上（按题目实际子问题数量增减）]

% ============================================================
% 二、问题分析
% ============================================================
\section{问题分析}

% 小节标题用"研究内容导向"（如"2.1 各催化剂组合指标与温度的关系分析"），
% 不用"问题1分析"这类标签式标题。可在此处放机理图、技术路线图、流程图。

\subsection{[问题1研究内容导向的小节标题]}

[该问的核心难点、拟采用的方法路线及理由、与其他问的关联，200-400 字]

\subsection{[问题2研究内容导向的小节标题]}

[同上]

% ============================================================
% 三、模型假设
% ============================================================
\section{模型假设}

% 3-6 条，每条 1-2 句简明陈述；仅关键假设附一句理由。禁止每条都展开"适用条件/放宽影响"。
\begin{enumerate}[label=\arabic*., leftmargin=2em]
    \item [假设内容，1-2 句]；
    \item [假设内容]；
    \item [假设内容]。
\end{enumerate}

% ============================================================
% 四、符号说明
% ============================================================
\section{符号说明}

% 默认单表三列（符号/说明/单位），与术语表.json 一致。
% 条件触发：特征体系庞大时，可分为"目标变量/特征"与"预处理参数"两张表。
% 说明须准确具体（如"单位 X 的 Y 成本"），禁止笼统"价格""系数"。
\begin{table}[H]
\centering
\caption{符号说明}
\begin{tabular}{ccc}
\toprule
符号 & 说明 & 单位 \\
\midrule
$x_1$ & [准确具体的含义] & [单位，无量纲则填 /] \\
$x_2$ & [含义] & [单位] \\
\bottomrule
\end{tabular}
\end{table}

% ============================================================
% 五、模型的建立与求解
% ============================================================
% 主体按问题组织：每个问题内部依次完成 建立→求解→验证→结果分析。
% 推导（公式由来、参数定义、模型性质）必须内联在正文对应位置，
% 不得后置到附录；附录仅放代码、支撑材料列表与确实过长的补充推导。
\section{模型的建立与求解}

\subsection{[问题1研究内容导向的小节标题]}

\subsubsection{数据探索与初步判断}
% C 题必写：分布观察/异常识别/结构发现 → 初步规律 → 处理决策，附图表；
% 无数据子问题可改为"问题结构分析"并说明理由。

[数据观察 → 初步规律 → 异常点处理，附图表]

\subsubsection{模型建立}

[模型假设如何落到数学形式。公式在此内联推导：先给形式，再逐项解释参数含义。]
\begin{equation}
y = f(x_1, x_2, \ldots; \theta)
\end{equation}
其中 $\theta$ 为[参数含义]；$x_i$ 为[含义]。

[理论性质（条件触发）：统计类→前提检验/拟合优度；机器学习类→模型结构与数据假设匹配性]

\subsubsection{模型求解}

[算法选择：≥2 种候选对比 + 选择理由（≥1 段）。]

[求解过程与结果。]

\subsubsection{结果验证与结果分析}

[样本外性能验证：验证设计（holdout/K 折/时序滚动，与数据结构匹配）结果，
均值±标准差；禁止单次划分结论。]

[模型验证：过拟合自查（train/test 双口径）+ 模型诊断（回归/时序→残差诊断；
检验类→前提检验）。蒙特卡洛仿真条件触发：含随机性/不确定性/稳态假设时，
报告仿真次数、收敛情况、与解析结果偏差、偏差 >5\% 的原因分析。
样本外性能验证与模型验证是两个独立维度，不可相互替代。]

[结果表格 + 图表引用 + ≥150 字解读（趋势/峰值/拐点 + 模型对应关系 + 实际意义）]

\subsection{[问题2研究内容导向的小节标题]}

[同上结构]

% ============================================================
% 六、灵敏度与误差分析
% ============================================================
\section{灵敏度与误差分析}

\subsection{灵敏度分析}

[数据扰动灵敏度（三选一以上）：输入特征 ±5\%/±10\%/±20\% 扰动 /
bootstrap 重抽样 / 多种子划分随机性，报告结论的稳定性。
扰动范围必须与代码实现一致。]
\begin{figure}[H]
\centering
\includegraphics[width=0.8\textwidth]{fig_sensitivity.pdf}
\caption{数据扰动灵敏度分析图}
\label{fig:sensitivity}
\end{figure}

% 条件触发小节（按题型保留，不适用则整节删除）：
\subsection{临界点分析（条件触发：结论随参数/阈值切换的问题）}

[参数/阈值在何值时结论发生转变，临界点图或表格。
所有临界点数值必须有 solution.json 或图表数据支撑，禁止凭直觉。]

\subsection{特征重要性与模型诊断（条件触发：机器学习/统计建模类问题）}

[特征重要性排序与业务解释；模型诊断结果（残差/学习曲线/混淆矩阵）汇总讨论。]

\subsection{误差来源与鲁棒性}

[误差来源 + 量化 + 影响 + 可接受性判断；鲁棒性验证。]

% ============================================================
% 七、模型评价与推广
% ============================================================
\section{模型评价与推广}

\subsection{优势}

[模型优势]

\subsection{局限性}

[模型局限性]

\subsection{改进方向与推广}

[改进方向；可推广到哪些场景]

\subsection{题目设计合理性讨论}

[批判性讨论题目设计：数据口径是否清晰？指标定义是否完整？样本代表性如何？假设是否需放松？]

% ============================================================
% 参考文献
% ============================================================
\begin{thebibliography}{99}
\bibitem{ref1} 作者. 标题[J]. 期刊, 年, 卷(期): 页码.
\bibitem{ref2} 作者. 标题[M]. 出版社, 年.
% ≥10 篇，GB/T 7714 格式，可查证可访问
\end{thebibliography}

% ============================================================
% 附录（官方：附录页数不限；必须含完整可运行源程序与支撑材料文件列表）
% ============================================================
\appendix

\section{支撑材料文件列表}

[列出随论文提交的支撑材料文件及用途；无则注明"本论文没有支撑材料"。]

\section{源程序代码}

[建模所用到的全部完整、可运行的源程序代码（官方硬约束：
缺可运行程序或程序运行结果与论文不符，可取消评奖资格）。
附录代码必须与 project/ 中实际代码一致，禁止"示意性改写"。]

\begin{lstlisting}[caption={问题1求解代码}]
# Python 代码（与实际求解代码一致）
import numpy as np

def solve_problem1():
    # 实际可运行代码
    pass
\end{lstlisting}

\section{补充推导（仅限正文确实放不下时）}

[正文模型建立处已给出完整推导链条；此处仅收录过长的中间代数步骤。
若使用本节，正文对应位置必须有明确引用（"详细中间步骤见附录 C"）。
禁止"详细推导较为复杂，从略"类回避性表述。]

\end{document}
